\begin{theorem}[Monotone stopping-ratio brackets]\label{thm:genie-ratio}
For every admissible quadruple of priors, with information ratios
interpreted in the extended nonnegative reals,
\begin{equation}\label{eq:genie-bracket}
 K_N\uparrow K_\infty\le K_N^+,
 \qquad K_{N+1}^+\le K_N^+.
\end{equation}
If, in addition, $C_v<\infty$, then
\begin{equation}\label{eq:genie-error}
 0\le K_N^+-K_N\le e_N^v,\qquad e_N^v\longrightarrow0.
\end{equation}
Consequently $K_N^+\downarrow K_\infty$ under this finite-divergence
hypothesis. Without it, the monotone brackets remain valid, but
pointwise convergence of the upper bracket need not hold.
\end{theorem}

\begin{remark}[Pointwise convergence can fail at infinite latent divergence]
Let $\mu(dd)=2\,dd$ on $(0,1/2]$, $\rho=1/4$, and take
$q_L=0$, $q_H=q_w=1/4$, and $q_v(d)=1/2-d$. All four priors are
admissible. Here
\[
 C_v=\frac18\int_0^{1/2}\frac{dd}{1/4-d^2}=\infty.
\]
The policy which always reaches depth $N$ and accepts the latent reveal
has infinite information and work $N$, so $K_N^+=\infty$ for every
finite $N$.

Nevertheless $K_\infty\le3/4$. To verify this for every capped policy,
parameterize the latent state by its Bernoulli mean $p$. The variance
reference has density $2(1-p)$ on $(0,1)$, the work reference has density
$3/2$ below $1/2$ and $1/2$ above $1/2$, and the signed endpoint
difference has density $-1/2$ below $1/2$ and $1/2$ above $1/2$.
For an ordered prefix $s$ of length $n\ge1$ with $k$ heads, put
$u_s=\int_0^1p^k(1-p)^{n-k}\,dp$. Then
\[
 |p_H(s)-p_L(s)|\le u_s/2,\qquad
 p_v(s)=\frac{2(n-k+1)}{n+2}u_s,
\]
and hence
\[
 r_s\le\frac{n+2}{8(n-k+1)}u_s\le\frac{3n}{8}u_s.
\]
Sum over stopped prefixes and the independent seed. If $U$ denotes the
auxiliary uniform law of $p$, the information of any capped policy is
at most $(3/8)\E_U T$. The work-reference density dominates half the
uniform density, so this is at most $(3/4)\E_{q_w}T$. Taking suprema
gives the claimed bound. The auxiliary law $U$ need not be admissible;
it is used only for this domination calculation.

This example does not contradict the optimized convergence in the next
theorem, whose proof first regularizes the variance reference.
\end{remark}

